%!TEX root=../../note.tex
\subsection{Summation and Product Notation}
Writing $1+2+3+4+5+6+7+8+9+10$ is manageable, but writing out every term
through $1000$ is not. Summation notation specifies both the terms and their
range, instead of relying on dots in $1+2+\cdots+1000$.

\begin{definition}[Summation notation]
    For integers $a\leq b$ and terms $t_a,\ldots,t_b$, write
    \begin{equation*}
        \sum_{i=a}^{b}t_i=t_a+t_{a+1}+\cdots+t_b.
    \end{equation*}
    The index $i$ runs through every integer from $a$ to $b$.
\end{definition}

\begin{example}
    \begin{equation*}
        1+2+\cdots+100=\sum_{i=1}^{100}i,
        \qquad
        1+8+27+64+125=\sum_{i=1}^{5}i^3.
    \end{equation*}
    In the second sum, the term indexed by $i$ is $i^3$.
\end{example}

\subsubsection{Properties of Summations}
For terms $x_i,y_i$, a constant $c$, and integer bounds $a\leq b$ and
$m\leq n$, the following identities let us rewrite sums without changing
their terms:
\begin{enumerate}
    \item Multiplying by a constant:
    \begin{equation*}
        c\sum_{i=a}^{b}x_i=\sum_{i=a}^{b}cx_i.
    \end{equation*}
    \item Adding or subtracting sums with the same bounds:
    \begin{equation*}
        \sum_{i=a}^{b}x_i\pm\sum_{i=a}^{b}y_i
        =\sum_{i=a}^{b}(x_i\pm y_i).
    \end{equation*}
    \item Shifting the index by an integer $r$:
    \begin{equation*}
        \sum_{i=m}^{n}x_i=\sum_{j=m+r}^{n+r}x_{j-r}.
    \end{equation*}
    The new index $j=i+r$ names the same terms.
    \item Splitting a sum at $n$, where $2\leq n\leq k$:
    \begin{equation*}
        \sum_{i=1}^{n-1}x_i+\sum_{i=n}^{k}x_i=\sum_{i=1}^{k}x_i.
    \end{equation*}
\end{enumerate}

\mdfsetup{nobreak=true}
\begin{example}
    Shifting by $r=1$ and splitting off the last term,
    \begin{equation*}
        \sum_{i=0}^{9}(i+1)=\sum_{j=1}^{10}j=\sum_{j=1}^{9}j+10.
    \end{equation*}
\end{example}
\mdfsetup{nobreak=false}

\subsubsection{Product Notation}
Products use the same indexing convention, with multiplication in place of
addition. For $n\geq 1$,
\begin{equation*}
    \prod_{i=1}^{n}i=1\cdot2\cdot3\cdots n=n!.
\end{equation*}

\subsubsection{Sums and Products with One or No Terms}
When the lower and upper bounds agree, the index takes a single value, so
the sum or product consists of one term:
\begin{equation*}
    \sum_{i=m}^{m}x_i=x_m \qquad\text{and}\qquad \prod_{i=m}^{m}x_i=x_m.
\end{equation*}

When the upper bound is smaller than the lower bound, no integer lies in the
range, and the sum or product is \emph{empty}. By convention an empty sum is
$0$ and an empty product is $1$. For example,
% TODO(check): the term of the empty sum was not recorded in lecture; $x_j$ is a placeholder.
\begin{equation*}
    \sum_{j=5}^{4}x_j=0 \qquad\text{and}\qquad \prod_{i=10}^{5}i^2=1.
\end{equation*}
\begin{remark}
    These values are the identities for $+$ and $\cdot$: adding $0$ or
    multiplying by $1$ changes nothing. With this convention the splitting
    identity above also holds for $n=1$, where the first piece is empty:
    $\sum_{i=1}^{0}x_i+\sum_{i=1}^{k}x_i=0+\sum_{i=1}^{k}x_i$.
    It also gives $0!=\prod_{i=1}^{0}i=1$.
\end{remark}

\subsubsection*{In practice}
\begin{itemize}
    \item \textbf{Reading a sum.} Identify the index, its first and last
    value, and the term as a function of the index; count the terms as
    $b-a+1$ (zero when $b<a$).
    \item \textbf{Rewriting a sum.} To compare two sums, shift one index so
    both have the same bounds, or split off the first or last term.
\end{itemize}